\lysh{36677}{4}

\begin{alist}
\item Από τα δεδομένα προκύπτει ότι:
\begin{rlist}
\item το πρόγραμμα Α περιγράφεται από μια γεωμετρική πρόοδο με $\alpha_1 = 1, \alpha_2 = 2, \alpha_3 = 4$ και $\lambda = 2$. Ισχύει επομένως ότι:
\[\alpha_\nu = \alpha_1 \cdot \lambda^{\nu-1} = 1 \cdot 2^{\nu-1} = 2^{\nu-1}\]
\item το πρόγραμμα Β περιγράφεται από μια αριθμητική πρόοδο με $\beta_1 = 100, \beta_2 = 110, \beta_3 = 120$ και $\omega = 10$. Ισχύει επομένως ότι:
\[\beta_\nu = \beta_1 + (\nu - 1) \cdot \omega = 100 + (\nu - 1) \cdot 10 = 100 + 10\nu - 10 = 10\nu + 90\]
\item Το ποσό που θα υπάρχει μετά από $\nu$ μήνες σύμφωνα με το πρόγραμμα Α θα είναι:
\[A_\nu = \alpha_1 \cdot \dfrac{\lambda^\nu - 1}{\lambda - 1} = 1 \cdot \dfrac{2^\nu - 1}{2 - 1} = 2^\nu - 1\]
\item Το ποσό που θα υπάρχει μετά από $\nu$ μήνες σύμφωνα με το πρόγραμμα Β θα είναι:
\[B_\nu = \dfrac{(\beta_1 + \beta_\nu) \cdot \nu}{2} = \dfrac{(100 + 10\nu + 90) \cdot \nu}{2} = \dfrac{(10\nu + 190) \cdot \nu}{2} = \dfrac{10\nu^2 + 190\nu}{2} = 5\nu^2 + 95\nu\]
\end{rlist}

\item
\begin{rlist}
\item Το ποσό που θα υπάρχει, σύμφωνα με το πρόγραμμα Α, μετά από $6$ μήνες, είναι:
$A_6 = 2^6 - 1 = 63\text{ \euro}$

Το ποσό που θα υπάρχει, σύμφωνα με το πρόγραμμα Β, μετά από $6$ μήνες, είναι:
\[B_6 = 5 \cdot 6^2 + 95 \cdot 6 = 180 + 570 = 750\text{ \euro}\]

\item Το ποσό που θα υπάρχει, σύμφωνα με το πρόγραμμα Α, μετά από $12$ μήνες, είναι:
$A_{12} = 2^{12} - 1 = 4095\text{ \euro}$

Το ποσό που θα υπάρχει, σύμφωνα με το πρόγραμμα Β, μετά από $12$ μήνες, είναι:
\[B_{12} = 5 \cdot 12^2 + 95 \cdot 12 = 720 + 1140 = 1860\text{ \euro}\]

Επομένως, ακολουθώντας το πρόγραμμα Α, θα έχει συγκεντρώσει μεγαλύτερο ποσό.
\end{rlist}
\end{alist}

